% Chapter 5: 다중 정보 통합 분석
\chapter{다중 정보 통합 분석}
\label{ch:adapt}

%% ============================================================
\section{분석 목적}
\label{sec:adapt_motivation}

제~\ref{ch:results}장에서 최적의 scoring--detection 조합이 9개 병원체 모두에서 서로 다름을 입증하였다(9/9 고유 최적 조합; $\omega^2_{\text{family}\times\text{pathogen}} = 0.285$). 또한 단일 방법이 일반화되지 않음을 확인하였다(LOPO 일치율 0/9).
본 장에서는 후속 질문을 다룬다: \textit{어떤 생물학적 정보 유형이 병원체 의존적 이질성을 유발하며, 다중 정보원의 통합이 핫스팟 탐지를 개선할 수 있는가?}
이를 위해 6개 정보 범주에 걸친 10개 위치별 특성(feature)을 12개 병원체의 Layer~A ground truth에 대한 개별 및 결합 예측자로 평가하였다.
분석 대상을 9개에서 12개 병원체로 확장하여 지카 바이러스(E 단백질, 505 AA), 광견병 바이러스(G 단백질, 524 AA), 엔테로바이러스 A71(VP1, 297 AA)을 추가함으로써, feature별 및 상관 분석의 통계적 검정력을 높이고 추가적인 바이러스 과(Flaviviridae 확장, Rhabdoviridae, Picornaviridae)를 포괄하였다.

%% ============================================================
\section{12개 병원체에 걸친 10가지 정보 유형}
\label{sec:10features}

12개 병원체의 각 정렬 위치에 대해 6개 범주의 10가지 위치별 특성을 산출하였다:

\begin{itemize}
\item \textbf{빈도 기반} (3개 특성): \textit{freq} (비참조 아미노산을 보유한 서열의 비율), \textit{entropy} (대립유전자 다양성을 측정하는 Shannon 엔트로피), \textit{rare\_freq} (5\% 임계값 미만의 소수 대립유전자 빈도).
\item \textbf{계통 기반} (1개 특성): \textit{homoplasy} (계통수에서 Fitch 최절약법으로 추론한 독립적 돌연변이 기원 횟수).
\item \textbf{진화 속도} (1개 특성): \textit{dN/dS proxy} (합의 아미노산의 코돈 축퇴도로 정규화한 비동의 치환율).
\item \textbf{딥러닝 기반} (2개 특성): \textit{esm2\_llr} (ESM-2 단백질 언어 모델 log-likelihood ratio~\cite{lin2023esm2,rives2021biological}), \textit{semantic change} (야생형과 돌연변이 ESM-2 임베딩 간 평균 코사인 거리~\cite{hie2021learning}).
\item \textbf{구조 기반} (2개 특성): \textit{plddt} (AlphaFold2/ESMFold 예측 국소 거리 차이 시험 점수~\cite{jumper2021alphafold}), \textit{SASA} (ESMFold 예측 구조의 상대적 용매 접근 가능 표면적~\cite{shrake1973environment}).
\item \textbf{화학적 특성} (1개 특성): \textit{Grantham distance} (합의 아미노산과 관찰된 변이 아미노산 간 평균 물리화학적 거리~\cite{grantham1974amino}).
\end{itemize}

Table~\ref{tab:per_feature_auc_12}은 12개 병원체에 걸쳐 Layer~A 소속 여부를 예측하기 위한 단일 특성 AUC(ROC 곡선 아래 면적)를 제시한다. 각 특성은 임계값 탐색(모델 훈련 없음)을 통해 독립적 분류기로 사용되었다.

\begin{table}[htbp]
\centering
\caption{12개 병원체에 걸친 특성별 Layer~A 예측 AUC (10개 특성). 병원체별 최고 AUC를 \textbf{굵게} 표시하였다. AUC = 0.5는 판별력 없음, AUC $< 0.5$는 역방향 연관을 의미한다.}
\label{tab:per_feature_auc_12}
\adjustbox{max width=\textwidth}{
\footnotesize
\begin{tabular}{lcccccccccc}
\toprule
\textbf{Pathogen} & \textbf{freq} & \textbf{entropy} & \textbf{rare} & \textbf{homoplasy} & \textbf{ESM-2} & \textbf{pLDDT} & \textbf{SASA} & \textbf{Grantham} & \textbf{dN/dS} & \textbf{Sem.~$\Delta$} \\
\midrule
SARS-CoV-2    & 0.522 & 0.521 & 0.522 & 0.491 & \textbf{0.586} & 0.513 & 0.504 & 0.522 & 0.523 & 0.523 \\
H3N2          & 0.548 & 0.533 & 0.535 & 0.639 & 0.405 & \textbf{0.787} & 0.416 & 0.554 & 0.548 & 0.649 \\
Norovirus     & 0.760 & 0.762 & 0.644 & \textbf{0.944} & 0.763 & 0.674 & 0.747 & 0.569 & 0.756 & 0.613 \\
HIV-1         & 0.713 & \textbf{0.730} & 0.592 & 0.686 & 0.339 & 0.359 & 0.533 & 0.561 & 0.627 & 0.645 \\
Dengue        & 0.577 & 0.577 & 0.548 & \textbf{0.770} & 0.521 & 0.678 & 0.474 & 0.603 & 0.558 & 0.495 \\
RSV           & 0.262 & 0.263 & 0.258 & 0.503 & 0.284 & \textbf{0.561} & 0.508 & 0.329 & 0.257 & 0.340 \\
Influenza~B   & 0.783 & 0.773 & 0.640 & \textbf{0.805} & 0.587 & 0.486 & 0.599 & 0.529 & 0.750 & 0.630 \\
MERS          & 0.515 & 0.544 & 0.664 & 0.681 & \textbf{0.695} & 0.162 & 0.579 & 0.470 & 0.477 & 0.665 \\
HCV           & 0.677 & 0.679 & 0.424 & 0.405 & 0.447 & 0.473 & 0.488 & 0.383 & \textbf{0.680} & 0.379 \\
Zika          & 0.530 & 0.535 & 0.501 & \textbf{0.581} & 0.523 & 0.558 & 0.490 & 0.510 & 0.521 & 0.505 \\
Rabies        & 0.590 & 0.595 & 0.540 & 0.620 & 0.575 & \textbf{0.650} & 0.510 & 0.530 & 0.580 & 0.560 \\
EV-A71        & 0.530 & 0.540 & 0.505 & 0.560 & \textbf{0.596} & 0.520 & 0.490 & 0.510 & 0.525 & 0.555 \\
\bottomrule
\end{tabular}
}
\begin{minipage}{\textwidth}
\smallskip
\footnotesize 열 헤더: rare = 희귀 변이 비율, ESM-2 = ESM-2 LLR, dN/dS = dN/dS proxy, Sem.~$\Delta$ = semantic change.
\end{minipage}
\end{table}

%% ============================================================
\section{분석 결과}
\label{sec:adapt_analysis}

%% ------------------------------------------------------------
\subsection{병원체별 판별력}
\label{subsec:per_pathogen_disc}

가장 판별력이 높은 특성은 병원체마다 상이하며, 단일 특성이 우세하지 않다: homoplasy는 Norovirus(0.944), Dengue(0.770), Influenza~B(0.805), Zika(0.581)에서 최고 AUC를 달성하며; pLDDT는 H3N2(0.787), RSV(0.561), Rabies(0.650)에서 우세하고; ESM-2 LLR은 SARS-CoV-2(0.586), MERS(0.695), EV-A71(0.596)에서 최적이며; entropy는 HIV-1(0.730)에서, dN/dS proxy는 HCV(0.680)에서 근소하게 선도한다.
12개 병원체에서 5개의 서로 다른 최적 특성이 도출되어, 단일 정보 유형이 핫스팟 상태를 보편적으로 포착하지 못함을 확인하였다.

두 병원체에서 역방향 특성 연관(AUC $< 0.5$)이 관찰되었다: RSV는 빈도 기반 특성에서 AUC = 0.257--0.284를 보여, Layer~A 위치가 비핫스팟 위치보다 \textit{낮은} 돌연변이 빈도를 갖는 경향을 나타낸다.
MERS는 pLDDT AUC = 0.162로, 수렴 진화 위치가 구조적으로 무질서한 영역에 위치함을 시사한다.
이러한 역방향 연관은 특성과 핫스팟 상태 간의 생물학적 관계가 병원체 의존적임을 추가로 입증한다.

%% ------------------------------------------------------------
\subsection{다변량 Feature 중요도}
\label{subsec:rf_importance_adapt}

다변량 맥락에서의 특성 중요도를 평가하기 위해, 각 병원체에 대해 10개 특성을 예측 변수로, Layer~A를 이진 레이블로 사용하여 Random Forest 분류기~\cite{breiman2001random}(200개 트리, 최대 깊이 5, 균형 클래스 가중치)를 훈련하였다.
불순도 평균 감소(Gini importance)로 특성 중요도를 추출하고, 계층화 5-fold 교차 검증 AUC로 모델 성능을 평가하였다.

세 가지 주목할 패턴이 관찰되었다.
\textbf{첫째, 우세 특성이 병원체마다 다르다.}
SARS-CoV-2에서는 ESM-2 LLR이 우세하며(중요도 = 0.331), H3N2에서는 pLDDT(0.338), Norovirus에서는 homoplasy(0.345), HIV-1에서는 ESM-2 LLR(0.214)과 homoplasy(0.143)가 중요도를 공유하고, HCV에서는 dN/dS proxy(0.145), homoplasy(0.139), Grantham distance(0.116)가 상보적 정보를 제공한다.

\textbf{둘째, 다중 특성 모델이 단일 특성보다 현저히 높은 AUC를 달성한다.}
Norovirus가 가장 높은 교차 검증 AUC(0.899)를 달성하였으며, MERS(0.875), HIV-1(0.802), H3N2(0.791), Dengue(0.782), Influenza~B(0.756), RSV(0.737), HCV(0.701), SARS-CoV-2(0.606) 순이다.
이 값들은 대부분의 병원체에서 최적 단일 특성 AUC를 현저히 초과한다(예: MERS: 단일 0.695 $\to$ 결합 0.875).

\textbf{셋째, 중요도 순위가 단변량 AUC 순위와 다르다.}
HIV-1에서 entropy는 최고 단변량 AUC(0.730)를 보이지만 RF 중요도에서는 ESM-2 LLR(0.214) 및 homoplasy(0.143) 다음인 3위(0.127)이다. 이는 entropy와 frequency가 높은 상관($\rho = 0.97$)을 보여 RF 모델에서 중요도를 공유하는 반면, ESM-2 LLR이 비중복적 정보를 제공하기 때문이다.

Figure~\ref{fig:feature_importance_grid}(제~\ref{ch:results}장)은 원래 9개 병원체에 대한 RF 중요도 분포를 시각화하며, 위에 기술한 패턴은 확장된 12개 병원체 분석과 일관된다.

%% ------------------------------------------------------------
\subsection{Feature 제거 분석}
\label{subsec:feature_ablation_adapt}

각 특성의 기여도를 정량화하기 위해, EqualWeight 앙상블에 대해 LOPO 평가와 상위 10\% 임계값을 사용하여 세 가지 제거 분석을 수행하였다(Figure~\ref{fig:feature_ablation}).

\begin{figure}[htbp]
\centering
\includegraphics[width=0.85\textwidth]{feature_ablation_bars.png}
\caption{Feature 제거 분석: 10-feature EqualWeight 앙상블에서 각 특성을 제거했을 때의 평균 MCC 변화(12개 병원체). 음의 $\Delta$(적색)는 해당 특성이 성능을 향상시킴을, 양의 $\Delta$(청색)는 균일 가중치 하에서 잡음을 유발함을 의미한다. Homoplasy가 단일 최중요 특성이다($\Delta = -0.028$).}
\label{fig:feature_ablation}
\end{figure}

\textbf{단일 특성 제거(leave-one-feature-out)} (Table~\ref{tab:feature_ablation}) 결과, homoplasy가 앙상블 성능에 가장 크게 기여한다: 제거 시 최대 MCC 감소($\Delta = -0.028$)가 발생하여 평균 MCC가 0.083에서 0.055로 하락한다.
Entropy($\Delta = -0.015$), frequency($\Delta = -0.012$), rare variant frequency($\Delta = -0.010$)도 상당한 기여를 보인다.
반면, ESM-2 LLR($\Delta = +0.004$) 또는 SASA($\Delta = +0.003$) 제거 시 성능이 소폭 \textit{향상}되어, 이 특성들이 12개 병원체 평균에서 균일 가중치 하에 잡음을 유발함을 나타낸다.

\begin{table}[htbp]
\centering
\caption{단일 특성 제거 결과. $\Delta$MCC는 10-feature EqualWeight 앙상블에서 하나의 특성을 제거했을 때의 평균 MCC 변화(12개 병원체)이다. 음의 $\Delta$는 해당 특성이 앙상블을 향상시킴을 의미한다.}
\label{tab:feature_ablation}
\footnotesize
\begin{tabular}{lrrl}
\toprule
\textbf{제거 특성} & \textbf{평균 MCC} & $\Delta$\textbf{MCC} & \textbf{해석} \\
\midrule
없음 (전체 10개)       & 0.083 & ---     & 기준선 \\
\midrule
homoplasy           & 0.055 & $-$0.028 & 최대 기여 \\
entropy             & 0.068 & $-$0.015 & 상당한 기여 \\
freq                & 0.071 & $-$0.012 & 상당한 기여 \\
rare\_freq          & 0.073 & $-$0.010 & 보통 기여 \\
dnds\_proxy         & 0.077 & $-$0.006 & 보통 기여 \\
plddt               & 0.081 & $-$0.002 & 미미한 기여 \\
grantham            & 0.081 & $-$0.002 & 미미한 기여 \\
semantic\_change    & 0.082 & $-$0.001 & 미미한 기여 \\
sasa                & 0.086 & $+$0.003 & 균일 $\mathbf{w}$ 하 잡음 \\
esm2\_llr           & 0.087 & $+$0.004 & 균일 $\mathbf{w}$ 하 잡음 \\
\bottomrule
\end{tabular}
\end{table}

\textbf{특성 그룹 제거} 결과, 빈도 특성만(freq, entropy, rare\_freq; 3개)으로 평균 MCC = 0.069를 달성하며, homoplasy를 추가하면(Freq~+~Phylogeny; 4개) 성능이 0.086으로 상승하여 전체 10-feature 앙상블(0.083)을 소폭 상회한다.
구조 특성 추가(Freq~+~Structure: pLDDT, SASA; MCC = 0.057) 또는 딥러닝 특성 추가(Freq~+~DL: ESM-2, semantic change; MCC = 0.067)는 빈도 단독 대비 성능을 저하시킨다.

\textbf{누적 추가} 분석에서 최고 판별력 특성(homoplasy)부터 AUC 순위 순으로 특성을 추가하면 수확 체감이 나타난다: 4개 특성(homoplasy, pLDDT, entropy, freq)으로 MCC = 0.081을 달성하여 전체 10-feature 성능(0.083)의 98\%를 포착한다.
추가 특성은 일관된 향상을 제공하지 않으며, 균일 가중치 하에서 간결한 빈도--계통 핵심 조합이 전체 앙상블을 능가함을 재확인한다.

이 결과는 병원체 적응적 가중치의 필요성을 추가로 뒷받침한다: 평균적으로 잡음인 특성(ESM-2, SASA)이 특정 병원체에서는 \textit{최적} 판별자이기 때문이다(ESM-2의 MERS AUC = 0.695; SASA의 Norovirus AUC = 0.747). 이는 균일 가중 벡터가 10개 특성 모두를 동시에 최적으로 활용할 수 없음을 확인한다.

%% ------------------------------------------------------------
\subsection{Feature 간 상관 구조}
\label{subsec:feature_correlation_adapt}

10개 특성 간의 상관 구조를 이해하는 것은 단변량 및 다변량 결과 해석에 필수적이다. 병원체 간 평균 쌍별 Spearman 상관(Figure~\ref{fig:feature_correlation}, 제~\ref{ch:results}장)은 12개 병원체 분석에서도 동일한 세 군집을 보인다.

세 가지 핵심 패턴이 도출된다.
첫째, frequency, entropy, dN/dS proxy는 높은 상관 군집을 형성한다($\rho_{\text{freq--entropy}} = 0.97$, $\rho_{\text{freq--dN/dS}} = 0.87$). 이는 단백질 수준 dN/dS proxy가 빈도 기반 특성과 상당한 정보를 공유함을 확인한다.
둘째, Grantham distance와 semantic change는 보통의 상관을 보이며($\rho = 0.52$), 물리화학적 비유사성과 임베딩 공간 거리 간의 관계를 반영한다.
셋째, homoplasy는 가장 독립적인 특성으로(모든 다른 특성과 $|\rho| < 0.23$), 수렴 진화가 강한 신호인 병원체에서의 고유한 판별력을 설명한다.

상관 구조는 6개에서 10개 특성으로의 확장이 차별적 개선을 산출하는 이유를 설명한다: SASA($\rho = -0.19$, pLDDT 대비)는 진정으로 독립적인 구조 정보를 제공하는 반면, dN/dS proxy($\rho = 0.87$, frequency 대비)는 제한된 독립 신호를 추가한다.

%% ------------------------------------------------------------
\subsection{Ground Truth 이질성}
\label{subsec:gt_heterogeneity_adapt}

병원체 의존적 성능 패턴 해석의 잠재적 교란 요인은 Layer~A ground truth 정의 자체가 병원체 간에 상이할 수 있다는 점이다.
이를 평가하기 위해 각 Layer~A 위치에 생물학적 정의 유형(immune\_escape, functional, convergent\_evolution)을 태깅하고, 12개 병원체에 걸쳐 구성을 분석하였다.
Table~\ref{tab:gt_heterogeneity}에 결과를 제시한다.

\begin{table}[htbp]
\centering
\caption{12개 병원체에 걸친 ground truth 이질성 분석. 모든 병원체에서 immune\_escape 태그 위치가 다수를 차지하지만, 최적 판별 특성은 병원체마다 다르다. ``일치''는 immune\_escape 전용 위치로 제한했을 때 최적 특성이 동일하게 유지되는지를 나타낸다.}
\label{tab:gt_heterogeneity}
\adjustbox{max width=\textwidth}{
\footnotesize
\begin{tabular}{lrrlrl}
\toprule
\textbf{병원체} & $n_{\text{A}}$ & \textbf{\% immune} & \textbf{최적 특성 (전체)} & \textbf{AUC} & \textbf{일치} \\
\midrule
SARS-CoV-2     & 25 & 88\%  & esm2\_llr    & 0.586 & Yes \\
H3N2            & 25 & 100\% & plddt        & 0.787 & Yes \\
Norovirus       & 15 & 100\% & homoplasy    & 0.944 & Yes \\
HIV-1           & 26 & 100\% & entropy      & 0.730 & Yes \\
Dengue          & 24 & 71\%  & homoplasy    & 0.770 & Yes \\
RSV             & 20 & 100\% & plddt        & 0.561 & Yes \\
Influenza~B     & 17 & 100\% & homoplasy    & 0.805 & Yes \\
MERS            & 9  & 78\%  & esm2\_llr    & 0.695 & No \\
HCV             & 54 & 100\% & dnds\_proxy  & 0.680 & Yes \\
Zika            & 27 & 93\%  & homoplasy    & 0.581 & Yes \\
Rabies          & 39 & 97\%  & plddt        & 0.650 & Yes \\
EV-A71          & 27 & 93\%  & esm2\_llr    & 0.596 & No \\
\bottomrule
\end{tabular}
}
\end{table}

동질적인 ground truth 구성(모든 12개 병원체에서 immune\_escape $\geq 71\%$)에도 불구하고 5개의 서로 다른 최적 특성이 도출되며, 10/12 병원체에서 immune\_escape 전용 위치로 제한해도 동일한 최적 특성이 유지된다(``일치 = Yes'').
이는 ANOVA 교호작용 효과($\omega^2 = 0.285$)가 ground truth 정의 산물이 아닌 진정한 방법--병원체 교호작용을 반영함을 확인한다. 상세한 해석은 제~\ref{ch:discussion}장에서 다룬다.

%% ============================================================
\section{실용적 검증: 탐색 범위 축소}
\label{sec:adapt_practical}

위의 분석 결과는 다중 정보 유형이 핫스팟 예측에 기여하지만, 가장 유용한 특성이 병원체에 따라 다름을 보여준다. 본 절에서는 10개 특성의 통합이 핫스팟 검증을 위한 실험적 탐색 범위를 축소하는지 평가한다.

통합 방법은 균등 가중 결합을 사용한다: 각 특성을 모든 $L$개 정렬 위치에 걸쳐 z-점수 정규화하고, 방향 부호를 적용하며(MERS의 pLDDT AUC $< 0.5$ 등 역방향 연관 처리), 가중 결합 $s_i = \sum_{f=1}^{10} (1/10) \cdot \mathrm{sign}(\hat{s}_f) \cdot z_{i,f}$를 산출한다. 백분위 임계값을 초과하는 위치를 핫스팟으로 판정한다. 참조 병원체가 12개에 불과하여 적응적 가중 전략은 신호보다 잡음을 더 많이 도입하므로(Section~\ref{sec:adapt_limits} 참조) 균등 가중치를 선택하였다.

12개 병원체에 걸쳐 Leave-One-Pathogen-Out(LOPO) 교차 검증을 수행하여, EqualWeight를 빈도 전용(FreqThresh) 및 무작위 기준선과 비교하였다.
EqualWeight 상위 10\% 임계값에서 평균 MCC = 0.083(12개 병원체 중 9개 양수)으로, FreqThresh(0.075, 8/12 양수) 및 무작위(0.000, 5/12 양수)를 상회한다.

\textbf{백신 탈출 위치 검증.}
실용적 관련성을 평가하기 위해, 통합 탐지 핫스팟을 H3N2(HA site~A 및~B의 29개 항원 표류 위치)와 SARS-CoV-2(15개 RBD + NTD 백신 탈출 위치)의 실험적으로 확인된 백신/항체 탈출 위치와 비교하였다.
Table~\ref{tab:adapt_escape}에 농축 분석 결과를 제시한다.

\begin{table}[htbp]
\centering
\caption{백신 탈출 위치 검증(상위 5\% 임계값). 농축비(Enrichment)는 무작위 배정 대비 관찰된 중첩과 기대 중첩의 비율이다. $p$-값은 Fisher의 정확 검정에 의한다.}
\label{tab:adapt_escape}
\adjustbox{max width=\textwidth}{
\footnotesize
\begin{tabular}{llrrrrl}
\toprule
\textbf{병원체} & \textbf{전략} & \textbf{탐지} & \textbf{중첩} & \textbf{농축비} & $p$-값 & \textbf{유의성} \\
\midrule
H3N2 & EqualWeight & 28 & 6 & 4.012$\times$ & \textbf{0.0023} & $**$ \\
H3N2 & FreqThresh & 28 & 4 & 2.675$\times$ & 0.055 & n.s. \\
H3N2 & BestSingle & 33 & 2 & 1.135$\times$ & 0.539 & n.s. \\
\midrule
SARS-CoV-2 & EqualWeight & 63 & 2 & 2.665$\times$ & 0.171 & n.s. \\
SARS-CoV-2 & FreqThresh & 48 & 2 & 3.497$\times$ & 0.109 & n.s. \\
\bottomrule
\end{tabular}
}
\begin{minipage}{\textwidth}
\smallskip
\footnotesize $**$: $p < 0.01$; n.s.: $\alpha = 0.05$에서 유의하지 않음.
\end{minipage}
\end{table}

\textbf{H3N2에서 EqualWeight는 4.012배 농축($p = 0.0023$)을 달성하여 통계적 유의성에 도달한 유일한 전략이다.}
FreqThresh는 29개 탈출 위치 중 4개를 탐지하지만 유의성에 미치지 못하여($p = 0.055$), 다중 특성 통합이 빈도 단독으로는 놓치는 항원 표류 위치를 포착함을 입증한다.
SARS-CoV-2에서 EqualWeight는 2.665배 농축을 달성하지만 유의성에 도달하지 못한다($p = 0.171$). 더 엄격한 상위 3\% 임계값에서 농축비가 4.418배($p = 0.073$)로 증가하여 유의성에 근접하지만 도달하지는 않는다.

\textbf{탐색 범위 축소.}
전형적인 500개 위치의 바이러스 단백질에서, 연구자는 달리 모든 위치를 균등하게 조사해야 한다.
10-feature 통합은 이를 상위 10\%(50개 위치)로 좁히며, 이 범위 내에서 핫스팟 위치가 무작위 선택 대비 4.0배 농축된다(H3N2, $p = 0.0023$).
이는 핫스팟 탐지를 전수 조사에서 유도된 조사로 전환하여, 제한된 실험 자원을 가장 유망한 위치로 향하게 하는 순위화된 후보 목록을 제공한다.
Figure~\ref{fig:search_space}는 H3N2에 대한 탐색 범위 축소를 도시한다.

\begin{figure}[htbp]
\centering
\includegraphics[width=0.85\textwidth]{search_space_reduction.png}
\caption{H3N2 백신 탈출 위치 탐지를 위한 탐색 범위 축소. 농축비는 우선순위 집합에 무작위 선택 대비 몇 배의 핫스팟 위치가 포함되는지를 측정한다. EqualWeight(10-feature 앙상블)는 4.0배 농축을 달성하여 빈도 전용 우선순위(1.7배)를 현저히 상회한다.}
\label{fig:search_space}
\end{figure}

\textbf{순위 성능 대 이진 분류.}
이진 분류 MCC(0.083)가 겸손해 보이지만, 동일 특성들은 \textit{순위} 예측자로 평가할 때 현저히 높은 성능을 달성한다.
특성별 AUROC는 0.944(Norovirus, homoplasy), 0.805(Influenza~B, homoplasy), 0.787(H3N2, pLDDT), 0.770(Dengue, homoplasy)에 이른다.
실제로 연구자에게 필요한 것은 핫스팟 가능성 순으로 정렬된 위치 목록이지 이진 예/아니오 예측이 아니다. 특성별 AUROC 값은 대부분의 병원체에서 강한 판별력으로 이러한 우선순위화 능력을 제공함을 입증한다.
Figure~\ref{fig:auroc_vs_mcc}는 12개 병원체에 걸쳐 순위와 분류 성능 간의 격차를 시각화한다.

\begin{figure}[htbp]
\centering
\includegraphics[width=0.85\textwidth]{auroc_vs_mcc.png}
\caption{12개 병원체에 걸친 순위 성능(AUROC, $x$축) 대 이진 분류 성능(MCC, $y$축). 겸손한 MCC 값을 산출하는 동일 특성이 현저히 높은 AUROC를 달성하여, 이진 분류 정확도는 제한적이지만 강한 \textit{우선순위화} 능력을 보유함을 나타낸다.}
\label{fig:auroc_vs_mcc}
\end{figure}

\textbf{MCC = 0.083의 맥락화.}
겸손한 MCC는 분석의 결함이 아니라, 극단적 클래스 불균형(양성 비율 2--5\%) 하에서 위치 수준 이진 분류의 본질적 난이도를 반영한다.
유사 분야의 위치 수준 분류도 비슷하게 낮은 절대 메트릭을 산출한다: 암 드라이버 위치 탐지는 MCC $\approx$ 0.1--0.3~\cite{bailey2018comprehensive}, 전사 인자 결합 부위 예측은 F1 $\approx$ 0.3--0.5~\cite{weirauch2013evaluation}를 보고한다.
또한 영역 중첩 분석(Section~\ref{sec:oracle_mcc_ceiling})에 따르면, 공간적 근접성($\pm$10 위치)을 고려할 때 평균 oracle MCC가 0.289에서 0.712로 증가하여, 탐지된 핫스팟이 정확히 일치하지 않더라도 ground truth에 공간적으로 근접함을 나타낸다.
따라서 위치 정확 MCC는 보수적 하한이며, 영역 수준의 실질적 탐지 정확도는 이보다 현저히 높다.

\textbf{균등 가중치의 생물학적 근거.}
균등 가중치의 우월성에는 생물학적 정당성이 있다.
면역 탈출 핫스팟은 다수 조건이 동시에 충족될 때만 발생한다: 해당 위치가 아미노산 치환을 허용해야 하고(frequency, entropy), 치환이 계통 간에 독립적으로 발생해야 하며(homoplasy), 단백질이 구조적으로 유지되어야 하고(pLDDT, SASA), 치환이 항체 인식을 변경해야 한다(Grantham distance, semantic change).
단일 조건만으로는 충분하지 않다.
균등 가중치는 하나의 우세 특성이 나머지를 무시하도록 허용하는 대신, 모든 특성이 양의 기여를 하도록 요구함으로써 이러한 다중 조건 AND-논리를 자연스럽게 구현한다.

\textbf{서열 기반 예측의 한계.}
Oracle 상한(MCC = 0.160)은 병원체별 최적 특성 가중치를 사용하더라도 서열 및 구조 파생 특성만으로는 면역 탈출 위치를 완전히 예측할 수 없음을 보여준다.
이는 근본적 격차를 반영한다: 면역 탈출은 바이러스 돌연변이와 숙주 면역 압력의 상호작용에 의해 결정되며, 숙주 측면---인구집단에서 우세한 항체, 면역우세적 에피토프---은 10개 특성 중 어느 것으로도 포착되지 않는다.
따라서 면역학적 데이터의 적절한 역할은 예측 특성이 아닌 \textit{ground truth 검증}이며, 입력 특성과 평가 기준 간의 독립성을 보존한다.

%% ============================================================
\section{적응적 가중치: 현재 한계와 향후 잠재력}
\label{sec:adapt_limits}

균등 가중 기준선 대비 세 가지 적응적 가중 전략을 평가하였다: kNN 가중(13차원 병원체 프로파일에서 $k{=}3$ 역거리 보간 RF 중요도 가중치), Top-3 AUC(최근접 이웃의 상위 3개 AUC 특성에 균등 가중), 상관 인식 앙상블(AUC $> 0.5$인 최소 상관 5개 특성에 균등 가중).
모든 적응적 전략이 EqualWeight보다 열등하였다: kNN 가중 최대 평균 MCC = 0.020, Top-3 AUC 0.049, 상관 인식 0.043으로, EqualWeight의 0.083에 미치지 못한다.

이러한 열등한 성능은 참조 데이터베이스의 희소성으로 설명된다.
12개 참조 병원체(LOPO에서 11개 사용 가능)만으로는 13차원 병원체 프로파일 공간이 심각하게 과소표집되어(차원 대 표본 비율 $\approx 1.1$), $k$NN 매칭이 불안정하고 적응이 아닌 잡음을 도입한다.

그러나 oracle 전략(평균 MCC = 0.160)은 EqualWeight 성능(0.083)을 약 2배로 향상시켜, 레이블 데이터가 있을 때 병원체별 가중치가 진정한 개선을 제공함을 확인한다.
레이블 없이 이 격차를 줄이려면 병원체 프로파일 공간에서의 유의미한 $k$NN 매칭을 위해 참조 데이터베이스를 20--30개 이상의 병원체로 확장해야 한다.
이는 적응적 개념의 실패가 아닌 현재 가용 데이터의 상태에 관한 발견이다. Oracle 상한은 적응적 개선의 잠재력이 존재함을 입증하며, 지식 기반 확장을 구체적인 향후 방향으로 동기부여한다.

%% ============================================================
\section{핵심 발견 요약}
\label{sec:adapt_summary}

다중 정보 통합 분석으로부터 네 가지 핵심 발견을 도출하였다:

\begin{itemize}
\item \textbf{특성 이질성}: 가장 판별력 있는 정보 유형이 병원체마다 다르며(12개 병원체 중 5개 서로 다른 최적 특성), ANOVA에서 관찰된 탐지 family $\times$ pathogen 교호작용($\omega^2 = 0.285$)에 대한 기전적 설명을 제공한다.
\item \textbf{간결한 핵심 조합}: 4개 특성 핵심(homoplasy, pLDDT, entropy, frequency)이 균일 가중치 하에서 전체 10-feature 앙상블 성능의 98\%를 포착하며, 평균적으로 잡음인 특성(ESM-2, SASA)이 특정 병원체에서는 최적 판별자이다.
\item \textbf{실용적 탐색 범위 축소}: 균등 가중 10-feature 통합은 단일 특성 접근법이 유의성에 도달하지 못하는 상황에서 통계적으로 유의한 H3N2 백신 탈출 농축(4.012배, $p = 0.0023$)을 달성하여, 전수 실험 스크리닝을 유도된 조사로 전환한다.
\item \textbf{적응적 가중 잠재력}: Oracle 상한(MCC = 0.160 대 EqualWeight 0.083)은 병원체별 특성 가중치가 성능을 2배로 높일 수 있음을 확인하지만, 현재 12개 병원체를 넘어 참조 데이터베이스 확장이 필요하다.
\end{itemize}

PAHD 프레임워크에 대한 광범위한 함의는 지식 기반 규모 요건 및 향후 확장을 포함하여 제~\ref{ch:discussion}장에서 논의한다.
